% ============================================================
%  preamble.tex —— 《正交宇宙物理学》共用导言区
%  编译： xelatex -> xelatex  (XeLaTeX + ctexbook)
% ============================================================
\documentclass[11pt, UTF8, scheme=chinese, fontset=windows]{ctexbook}

\usepackage{amsmath, amssymb, amsfonts, mathrsfs, bm}
\usepackage{graphicx}
\usepackage{geometry}
\geometry{a4paper, top=25mm, bottom=25mm, left=28mm, right=28mm}
\usepackage{booktabs}
\usepackage{array}
\usepackage{multirow}
\usepackage{longtable}
\usepackage{enumitem}
\usepackage{xcolor}
\usepackage{caption}
\usepackage[most]{tcolorbox}
\usepackage{fancyhdr}
\usepackage{hyperref}

\hypersetup{
  colorlinks=true,
  linkcolor=blue!50!black,
  urlcolor=teal!60!black,
  citecolor=green!45!black,
  bookmarksnumbered=true,
  pdftitle={正交宇宙物理学},
  pdfauthor={编译自 Greg Egan《Orthogonal》官方设定}
}

% ---- 图表 ----
\captionsetup{font=small, labelfont=bf, labelsep=quad}
\graphicspath{{images/}}

% 插图命令：\orthofig[宽度]{图文件名}{中文图注}
\newcommand{\orthofig}[3][0.62\textwidth]{%
  \begin{figure}[htbp]
    \centering
    \includegraphics[width=#1]{#2}
    \caption{#3}
  \end{figure}}

% 不浮动的小图（用于正文中就地的示意图）
\newcommand{\inlinefig}[3][0.5\textwidth]{%
  \begin{center}
    \includegraphics[width=#1]{#2}\\[2pt]
    {\small #3}
  \end{center}}

% ---- 自定义盒子 ----
\newtcolorbox{keybox}[1][]{
  colback=blue!4!white, colframe=blue!45!black,
  boxrule=0.6pt, arc=2pt, left=6pt, right=6pt, top=5pt, bottom=5pt,
  fonttitle=\bfseries, title=#1}

\newtcolorbox{notebox}[1][]{
  colback=orange!5!white, colframe=orange!65!black,
  boxrule=0.6pt, arc=2pt, left=6pt, right=6pt, top=5pt, bottom=5pt,
  fonttitle=\bfseries, title=#1}

% ---- 常用记号 ----
\newcommand{\Riem}{\mathrm{R}}
\newcommand{\Lor}{\mathrm{L}}
\newcommand{\dd}{\mathrm{d}}
\newcommand{\ii}{\mathrm{i}}
\newcommand{\ee}{\mathrm{e}}
\newcommand{\wms}{\omega_m}
\newcommand{\bA}{\mathbf{A}}
\newcommand{\bE}{\mathbf{E}}
\newcommand{\bB}{\mathbf{B}}
\newcommand{\bF}{\mathbf{F}}
\newcommand{\bj}{\mathbf{j}}
\newcommand{\bk}{\mathbf{k}}
\newcommand{\bp}{\mathbf{p}}
\newcommand{\bv}{\mathbf{v}}
\newcommand{\bu}{\mathbf{u}}
\newcommand{\bw}{\mathbf{w}}
\newcommand{\bx}{\mathbf{x}}
\newcommand{\be}{\mathbf{e}}
\newcommand{\bP}{\mathbf{P}}
\newcommand{\bq}{\mathbf{q}}
\newcommand{\bzero}{\mathbf{0}}
\newcommand{\half}{\tfrac{1}{2}}
\DeclareMathOperator{\Tr}{Tr}
\DeclareMathOperator{\diag}{diag}
\DeclareMathOperator{\sgn}{sgn}
\DeclareMathOperator{\Span}{span}
\DeclareMathOperator{\rank}{rank}

% ---- 页眉页脚 ----
\pagestyle{fancy}
\fancyhf{}
\fancyhead[LE,RO]{\small\thepage}
\fancyhead[LO]{\small\nouppercase{\rightmark}}
\fancyhead[RE]{\small\nouppercase{\leftmark}}
\renewcommand{\headrulewidth}{0.4pt}
